\section{Mathematical Definition and Properties}
\label{sec:definition}

\subsection{Definition}

For a $k$-district congressional plan with district-level Democratic vote shares
$v_1, v_2, \ldots, v_k$, partition the districts into Democratic-won ($\mathcal{D} =
\{i : v_i > 0.50\}$) and Republican-won ($\mathcal{R} = \{i : v_i < 0.50\}$). Let
$k_D = |\mathcal{D}|$ and $k_R = |\mathcal{R}|$ be the number of districts won
by each party.

Compute the mean Democratic vote share in each party's won districts:
\[
\bar{v}_D^+ = \frac{1}{k_D}\sum_{i \in \mathcal{D}} v_i, \qquad
\bar{v}_R^- = \frac{1}{k_R}\sum_{i \in \mathcal{R}} v_i.
\]
$\bar{v}_D^+$ is the mean of Democratic-won district vote shares (values above 0.50);
$\bar{v}_R^-$ is the mean of Republican-won district vote shares (values below 0.50).

Compute the angles:
\[
\theta_D = \arctan(2\bar{v}_D^+ - 1), \qquad \theta_R = \arctan(1 - 2\bar{v}_R^-).
\]
$\theta_D$ measures how far the center of Democratic-won districts is from 50\%:
if Democratic-won districts average $\bar{v}_D^+ = 0.65$, then $\theta_D = \arctan(0.30)
\approx 0.29$ radians. $\theta_R$ symmetrically measures how far the center of
Republican-won districts is from 50\%: if Republican-won districts average
$\bar{v}_R^- = 0.42$, then $\theta_R = \arctan(0.16) \approx 0.16$ radians.

Declination is:
\[
\delta = \frac{2(\theta_D - \theta_R)}{\pi}.
\]
The normalization by $\pi/2$ maps the range to $[-1, 1]$. Positive $\delta$ indicates
that Democratic-won districts are further from 50\% than Republican-won districts
($\theta_D > \theta_R$), i.e., Democratic voters are packed into high-margin
districts while Republican voters win by smaller margins---a configuration associated
with Republican structural advantage.

\subsection{Geometric Interpretation}

The geometric interpretation of Declination is most transparent on a plot of sorted
district vote shares. Sort all $k$ districts by Democratic vote share from lowest
to highest. Plot the vote shares on the vertical axis against district rank on the
horizontal axis. Draw a horizontal line at 50\%.

In a district with pure cracking of Democratic voters and pure packing of remaining
Democrats:
\begin{itemize}
\item Republican-won districts (left half of the plot) cluster near 45\%---close to
  the 50\% line, producing small $\theta_R$.
\item Democratic-won districts (right half) cluster near 75\%---far from the 50\%
  line, producing large $\theta_D$.
\end{itemize}
The declination $\delta = 2(\theta_D - \theta_R)/\pi$ is large positive.

In a compact, neutral redistricting plan, both parties' won districts average
similarly close to 50\%, producing $\theta_D \approx \theta_R$ and $\delta \approx 0$.

\subsection{Three Methodological Critiques}

\textbf{Critique 1: Undefined when one party wins all seats.}
If $k_R = 0$ (all districts won by Democrats) or $k_D = 0$ (all districts won by
Republicans), $\bar{v}_R^-$ or $\bar{v}_D^+$ is undefined, and Declination cannot
be computed. This is not merely a hypothetical concern: in small delegations in
wave elections, one party may sweep all districts, and in states with very lopsided
partisan geography (e.g., MA $k=9$ in a strong Democratic year), the metric is
undefined. Expert witnesses cannot rely on Declination in cases where such edge
conditions arise.

\textbf{Critique 2: Sensitive to the 50\% seat-win threshold.}
The partition of districts into $\mathcal{D}$ and $\mathcal{R}$ depends on which
districts are above and below 50\%. A district with $v_i = 0.501$ (barely Democratic)
contributes to $\bar{v}_D^+$; a district with $v_i = 0.499$ (barely Republican)
contributes to $\bar{v}_R^-$. Small changes in vote shares near 50\% can move
districts between categories, causing discontinuous changes in Declination. This
discontinuity is more severe than the analogous discontinuity in EG (where flipping
a district changes the wasted-vote count by a bounded amount).

\textbf{Critique 3: No established unconstitutionality threshold.}
No court has determined what value of Declination constitutes a constitutionally
problematic partisan advantage. \citet{warrington2018} proposes $|\delta| > 0.2$
as indicating substantial partisan advantage but acknowledges this is tentative.
Unlike EG's 8\% threshold (which was at least proposed explicitly), Declination
has no corresponding proposal with empirical calibration.
